% Nat-Elim + Catamorphisms + Recursion Zoo + HOF + Normalization (Ch 27-31)
% Lines 6844-7204 of main.tex
% This file is for agent editing — will be merged back

\chapter{The Nat-Eliminator Extraction}
\label{ch:nat-elim}

\section{The General Extraction Algorithm}

Given a Python AST for a recursive function, the Nat-eliminator
extraction identifies:
\begin{enumerate}
  \item The \textbf{base case}: the guard condition and return value
    when the parameter is ``small enough.''
  \item The \textbf{step operation}: the binary operation applied to
    the parameter and the recursive result.
  \item The \textbf{threshold}: the parameter value at which the
    base case triggers.
\end{enumerate}

The extraction is \emph{not algorithm-specific}: it works for any
recursive function of the form:
\begin{lstlisting}
def f(n):
    if n OP threshold: return base_val
    return BINOP(n, f(n - 1))
\end{lstlisting}
where \texttt{OP} is any comparison ($\leq$, $<$, $=$), \texttt{BINOP}
is any binary operation ($+$, $\times$, $-$, $//$, $\%$, $|$, $\&$,
$\oplus$), and either $n$ or $f(n-1)$ can be on either side of
\texttt{BINOP} (handled by commutativity when applicable).

\section{Correctness}

\begin{theorem}[Extraction Correctness]
If the Nat-eliminator extraction succeeds with parameters
$(\mathsf{base}, \mathsf{step}, \mathsf{threshold}, \op)$, then the
canonical fold $\fold(\op, \mathsf{threshold}, \mathsf{ival}(p) + 1,
\mathsf{IntObj}(\mathsf{base}))$ computes the same function as the
original recursive program.
\end{theorem}

\begin{proof}
The fold satisfies the same recurrence as the recursive program:
\begin{align}
  \fold(\op, k, k, \mathsf{base}) &= \mathsf{base}
    \quad \text{(base case: $i \geq \mathsf{stop}$)} \\
  \fold(\op, i, \mathsf{stop}, \mathsf{acc}) &=
    \mathsf{binop}(\op, \mathsf{IntObj}(i),
    \fold(\op, i+1, \mathsf{stop}, \mathsf{acc}))
    \quad \text{(step)}
\end{align}
This matches the recursive definition
$f(n) = \mathsf{binop}(\op, n, f(n-1))$ with $f(k) = \mathsf{base}$,
by the Nat-recursor uniqueness (Theorem D1).
\end{proof}


\chapter{Iteration Beyond Range: Collection Catamorphisms}
\label{ch:iteration-beyond}

The Nat-eliminator handles \texttt{for i in range(n)}.  Real programs
iterate over lists, dicts, strings, generators, and composed iterators.

\section{The List Catamorphism}

A list $[a,b,c]$ is the Pair-chain $\mathsf{Pair}(a, \mathsf{Pair}(b,
\mathsf{Pair}(c, \mathsf{NoneObj})))$.  Iteration over it is a
\emph{list fold} (catamorphism on the free monoid):

\begin{definition}[List Fold]
\label{def:list-fold}
For \texttt{for x in lst: acc = step(acc, x)}:
\[
  \mathsf{list\_fold}(\mathsf{lst}, \mathsf{acc}) = \begin{cases}
    \mathsf{acc} & \mathsf{is\_NoneObj}(\mathsf{lst}) \\
    \mathsf{list\_fold}(\mathsf{snd}(\mathsf{lst}),\;
      \mathsf{step}(\mathsf{acc}, \mathsf{fst}(\mathsf{lst})))
      & \mathsf{is\_Pair}(\mathsf{lst}) \\
    \mathsf{acc} & \text{otherwise}
  \end{cases}
\]
\end{definition}

\begin{lstlisting}[caption={List fold Z3 RecFunction}]
def compile_for_over_list(stmt, env, T):
    S = T.S; T._uid += 1; uid = T._uid
    target = stmt.target.id
    iter_term = compile_expr(stmt.iter, env)
    modified = find_modified(stmt.body)

    for var in modified:
        init = env.get(var)
        rfn = RecFunction(f'lfold_{var}_{uid}', S, S, S)
        lst_s = FreshConst(S, f'_lf_{uid}')
        acc_s = FreshConst(S, f'_la_{uid}')
        body_env = env.copy()
        body_env.put(target, S.fst(lst_s))
        body_env.put(var, acc_s)
        exec_stmts(stmt.body, body_env)
        new_val = body_env.get(var)
        RecAddDefinition(rfn, [lst_s, acc_s],
            If(S.is_NoneObj(lst_s), acc_s,
               If(S.is_Pair(lst_s),
                  rfn(S.snd(lst_s), new_val), acc_s)))
        env.put(var, rfn(iter_term, init))
\end{lstlisting}

\begin{theorem}[List Fold = Nat Fold on Length]
For a list of length $n$: $\mathsf{list\_fold}([a_1,\ldots,a_n], e)
= \fold(\mathsf{step}', 0, n, e)$ where $\mathsf{step}'(i, \mathsf{acc})
= \mathsf{step}(\mathsf{acc}, a_{i+1})$.  This connects list iteration
to range iteration via the D1 path.
\end{theorem}

\section{Dict Iteration}

\texttt{for k,v in d.items()} compiles as a list fold over the shared
function $\mathsf{py\_meth\_items}(d)$, which returns a Pair-chain of
key-value Pairs.

\begin{proposition}[Dict Iteration Order Independence]
If the fold operation is commutative and associative, the result is
independent of dict iteration order.  The commutative monoid paths
(F2, F3) provide the connecting path.
\end{proposition}

\section{String Iteration}

\texttt{for ch in s} is a Nat fold over the character index:
$\mathsf{str\_fold}(s, i, \mathsf{acc}) = \begin{cases}
\mathsf{acc} & i \geq \mathsf{Length}(s) \\
\mathsf{str\_fold}(s, i{+}1, \mathsf{step}(\mathsf{acc},
  \mathsf{StrObj}(\mathsf{SubStr}(s, i, 1)))) & \text{otherwise}
\end{cases}$

This integrates directly with the existing Nat fold machinery.

\section{Composed Iterators}

\texttt{zip}, \texttt{enumerate}, \texttt{map}, \texttt{filter}
transform iterators.  Each compiles to a shared function wrapping the
inner iterator:
\begin{lstlisting}[caption={Composed iterators}]
# for i,x in enumerate(xs): ...  => list_fold(py_enumerate(xs), acc)
# for a,b in zip(xs,ys): ...     => list_fold(py_zip(xs,ys), acc)
# sum(x*2 for x in xs)           => list_fold(py_map(h, xs), IntObj(0))
\end{lstlisting}

\section{Generator Consumption}

A generator yielding values is modeled as producing a Pair-chain.
Full consumption by \texttt{list()}, \texttt{for}, or \texttt{sum}
gives a list fold.  This is the D6 path: both generator and eager
versions produce the same Pair-chain.


\chapter{The Recursion Zoo: Beyond Nat Descent}
\label{ch:recursion-zoo}

\section{Tree Recursion}

\begin{definition}[Tree Catamorphism]
$f(\mathsf{tree}) = \begin{cases}
\mathsf{base}(\mathsf{tree}) & \mathsf{is\_leaf} \\
\mathsf{combine}(f(\mathsf{left}), f(\mathsf{right}), \mathsf{tree})
  & \text{otherwise}
\end{cases}$

This is the unique algebra morphism from the initial algebra of
binary trees.
\end{definition}

\begin{lstlisting}[caption={Tree recursion Z3 encoding}]
def compile_tree_rec(func, body, env, T):
    S = T.S; T._uid += 1; uid = T._uid
    param = func.args.args[0].arg
    sym = FreshConst(S, f'_tree_{uid}')
    rfn = RecFunction(f'tree_{func.name}_{uid}', S, S)
    rec_env = Env(T, func_name=func.name, is_rec=True)
    rec_env.put(param, sym)
    secs = extract_sections(body, rec_env)
    combined = secs[-1].term
    for s in reversed(secs[:-1]):
        combined = If(s.guard, s.term, combined)
    RecAddDefinition(rfn, [sym], combined)
    return rfn
\end{lstlisting}

\begin{theorem}[Tree Catamorphism Uniqueness]
Two functions with the same $(\mathsf{base}, \mathsf{combine})$ are
equal by the initial algebra universal property.  Stack-based traversal
with the same combine is connected by D9 (Stack$\leftrightarrow$Recursion).
\end{theorem}

\section{Multi-Parameter Recursion (2D)}

Edit distance, knapsack, LCS all recurse on two parameters:
$f(i,j) = \mathsf{step}(f(i{-}1,j), f(i,j{-}1), f(i{-}1,j{-}1))$.

\begin{lstlisting}[caption={2D RecFunction}]
f = RecFunction('f_2d', IntSort(), IntSort(), S)
i, j = Ints('i j')
RecAddDefinition(f, [i, j],
    If(i <= 0, base_row(j),
    If(j <= 0, base_col(i),
       step(i, j, f(i-1,j), f(i,j-1), f(i-1,j-1)))))
\end{lstlisting}

Both memoized and bottom-up DP satisfy the \emph{same} 2D recurrence.
By $\Nat \times \Nat$-eliminator uniqueness: D16 path.

\section{Mutual Recursion}

$f(n) = \ldots g(n{-}1)\ldots$, $g(n) = \ldots f(n{-}1)\ldots$.
Model as product recursor $h(n) = (f(n), g(n))$:

\begin{lstlisting}[caption={Mutual recursion as product}]
h = RecFunction('mutual', IntSort(), S)
RecAddDefinition(h, [n],
    If(n <= 0, S.Pair(base_f, base_g),
       S.Pair(step_f(n, S.fst(h(n-1)), S.snd(h(n-1))),
              step_g(n, S.fst(h(n-1)), S.snd(h(n-1))))))
\end{lstlisting}

\section{Well-Founded Recursion with Fuel}

General recursion on any well-founded relation, made terminating
with a fuel parameter:

\begin{lstlisting}[caption={Well-founded recursion}]
f = RecFunction('wf', S, IntSort(), S)
RecAddDefinition(f, [x, fuel],
    If(fuel <= 0, S.Bottom,
    If(is_base(x), base(x),
       step(x, f(smaller(x), fuel - 1)))))
\end{lstlisting}

Decidable for ground arguments (fuel bounds the depth).  Two
well-founded recursions with same $(\mathsf{base}, \mathsf{step},
\prec)$ connected by W-type eliminator uniqueness.


\chapter{Higher-Order Functions: Compilation and Normalization}
\label{ch:hof-practice}

\section{Compilation Rules}

\begin{lstlisting}[caption={HOF compilation}]
# map(lambda x: x*2, xs)
# => py_map_{hash("x*2")}(xs)          [shared symbol]

# filter(lambda x: x > 0, xs)
# => py_filter_{hash("x > 0")}(xs)     [shared symbol]

# sorted(xs, key=len)
# => py_sorted_by_{hash("len")}(xs)    [shared symbol]

# reduce(op, xs, init)
# => py_reduce_{hash(op)}(xs, init)    [shared symbol]
\end{lstlisting}

The key insight: two programs using the \emph{same} lambda body get
the \emph{same} shared symbol, regardless of whether it's written as
\texttt{map(f, xs)}, a comprehension \texttt{[f(x) for x in xs]}, or
an explicit loop.

\section{HOF Normalization Rules}

\begin{enumerate}
  \item $\mathsf{reduce}(\oplus, e, \mathsf{map}(h, xs)) \to
    \mathsf{reduce}(\lambda a,x.\, a \oplus h(x), e, xs)$
    \hfill (map--fold fusion)
  \item $\mathsf{map}(f, \mathsf{map}(g, xs)) \to
    \mathsf{map}(f \circ g, xs)$
    \hfill (map composition)
  \item $\mathsf{filter}(p, \mathsf{map}(h, xs)) \to
    \mathsf{map}(h, \mathsf{filter}(p \circ h, xs))$
    \hfill (filter--map swap, when applicable)
  \item $\mathsf{sorted}(\mathsf{list}(x)) \to \mathsf{sorted}(x)$
    \hfill (list absorption)
\end{enumerate}

\begin{example}[map+filter vs comprehension]
$f(xs) = \mathsf{list}(\mathsf{filter}_{h_1}(\mathsf{map}_{h_2}(xs)))$
vs.\ $g(xs) = \mathsf{comp}_{h_3}(xs)$.

Path: D4 identifies the comprehension with filter$\circ$map when
$h_3$ matches the composition of $h_1$ and $h_2$.
\end{example}


\chapter{Complete Normalization Traces}
\label{ch:traces}

\section{Trace 1: Factorial (Rec $\to$ Fold)}

\textbf{Input}: \texttt{def f(n): if n<=1: return 1; return n*f(n-1)}

\begin{enumerate}
  \item \textbf{Compilation}: $\den{f} = \mathsf{rec\_f}(p_0)$ where
    $\mathsf{rec\_f}(n) = \mathsf{If}(\mathsf{ival}(n) \leq 1,
    \mathsf{IntObj}(1), \mathsf{binop}(\mathsf{MUL}, n,
    \mathsf{rec\_f}(\mathsf{binop}(\mathsf{SUB}, n, \mathsf{IntObj}(1)))))$.
  \item \textbf{Phase 3 (Nat-eliminator)}: Detects base=$1$,
    threshold=$1$, op=$\mathsf{MUL}$.  Produces
    $\fold(\mathsf{MUL}, 1, \mathsf{ival}(p_0)+1, \mathsf{IntObj}(1))$.
  \item \textbf{Iterative version}: same fold via loop-to-fold.
  \item \textbf{Comparison}: identical Z3 terms.  Path: $\mathsf{refl}$.
\end{enumerate}

\section{Trace 2: Edit Distance (2D Recurrence)}

\textbf{Input A} (memoized): nested \texttt{dp(i,j)} with \texttt{@lru\_cache}.

\begin{enumerate}
  \item \textbf{Phase 1}: Inline nested \texttt{dp} (D2).  Ignore decorator.
  \item \textbf{Compilation}: 2D RecFunction $\mathsf{dp}(i,j) =
    \mathsf{If}(i=0, j, \mathsf{If}(j=0, i,
    \mathsf{binop}(\mathsf{MIN}, \ldots)))$.
  \item \textbf{Phase 2}: Canonicalize $\min$ argument order (L1).
  \item \textbf{Phase 5}: Shared symbols for subscript, len.
  \item \textbf{Input B} (DP table): same 2D RecFunction after
    compiling nested loops as 2D fold.
  \item \textbf{Comparison}: same recurrence, same bases.
    Path: $\mathsf{D16}$ (Memo$\leftrightarrow$DP).
\end{enumerate}

\section{Trace 3: all vs any (De Morgan)}

\textbf{Input A}: \texttt{all(pred(x) for x in seq)}.

\textbf{Input B}: \texttt{not any(not pred(x) for x in seq)}.

\begin{enumerate}
  \item \textbf{Phase 5}: A $\to \mathsf{py\_all}(\mathsf{comp}_{h_1}(p_0))$.
    B $\to \mathsf{unop}(\mathsf{NOT}, \mathsf{py\_any}(\mathsf{comp}_{h_2}(p_0)))$.
  \item $\mathsf{comp}_{h_2}$ has body \texttt{not pred(x)} =
    $\mathsf{map}(\mathsf{NOT}, \mathsf{comp}_{h_1})$.
  \item \textbf{Path B6} (De Morgan): $\mathsf{py\_all}(xs) \equiv
    \lnot\,\mathsf{py\_any}(\mathsf{map}(\lnot, xs))$.
  \item Certificate: 2 steps.
\end{enumerate}

\section{Trace 4: Sorted Unique}

\textbf{Input A}: \texttt{return sorted(set(xs))}.

\textbf{Input B}: manual dedup + insertion sort.

\begin{enumerate}
  \item \textbf{Phase 5}: A $\to \mathsf{py\_sorted}(\mathsf{py\_set}(p_0))$.
  \item B compiles to a list fold with insert operations.
  \item \textbf{Path}: requires axiom $\mathsf{py\_sorted}(\mathsf{py\_set}(x))$
    = canonical ``sorted unique elements''.  Both programs produce this
    canonical form.  Path: $\mathsf{I3} \cdot \mathsf{I5}$
    (set idempotence $+$ sorted idempotence).
\end{enumerate}


% ══════════════════════════════════════════════════════════
\part{Specification Checking and Bug Finding via CTTT}
% ══════════════════════════════════════════════════════════

